\documentclass[12pt]{article}

% Idioma e codificação
\usepackage[brazil]{babel}
\usepackage[utf8]{inputenc}
\usepackage[T1]{fontenc}

% Matemática e listas
\usepackage{amsmath, amssymb}
\usepackage{enumitem}

% Layout e cores
\usepackage{geometry}
\usepackage{graphicx}
\usepackage{xcolor}

% Cores institucionais
\definecolor{maincolor}{RGB}{0,100,50}
\definecolor{forestgreen}{RGB}{0,100,50}

% Cabeçalho e rodapé
\usepackage{fancyhdr}

% Blocos
\usepackage[most]{tcolorbox}
\newtcolorbox{block}[1]{%
  colback=green!5,
  colframe=maincolor,
  coltitle=white,
  title=#1,
  fonttitle=\bfseries,
  boxrule=1.5pt,
  arc=3mm,
  left=6pt,
  right=6pt,
  top=6pt,
  bottom=6pt
}

% Margens
\geometry{margin=2.5cm}

% Fancyhdr
\pagestyle{fancy}
\fancyhf{}

\renewcommand{\headrulewidth}{2pt}
\renewcommand{\footrulewidth}{2pt}

\renewcommand{\headrule}{\hbox to\headwidth{%
  \color{forestgreen}\leaders\hrule height \headrulewidth\hfill}}

\renewcommand{\footrule}{\hbox to\headwidth{%
  \color{forestgreen}\leaders\hrule height \footrulewidth\hfill}}

% Altura do cabeçalho (corrigida)
\setlength{\headheight}{52pt}
\addtolength{\topmargin}{-20pt}

% Cabeçalho
\fancyhead[L]{%
  \includegraphics[height=1.2cm]{Universidade_de_Rio_Verde.jpg}
}

\fancyhead[R]{%
  \raggedleft
  \textcolor{forestgreen}{%
    \small
    \textbf{Lista de Exercícios — Paradigmas de Programação}\\
    Me.~Sergio Souza Novak\\
    \textit{Programação Orientada a Objetos}
  }
}

% Rodapé
\fancyfoot[C]{\textcolor{forestgreen}{\thepage}}

% =========================
\begin{document}
% =========================

\vspace*{0.2cm}

\section*{Exercícios}

\begin{block}{Questão 1 — Conceitos Fundamentais}
Explique, com suas próprias palavras, o que é um \textbf{paradigma de programação}.
Por que diferentes paradigmas existem e qual a importância deles no desenvolvimento
de software?
\end{block}

\begin{block}{Questão 2 — Associação}
Associe cada paradigma à sua principal característica:

\begin{enumerate}[label=\alph*)]
  \item Paradigma Imperativo
  \item Paradigma Orientado a Objetos
  \item Paradigma Funcional
  \item Paradigma Lógico
  \item Paradigma Declarativo
\end{enumerate}

\begin{itemize}
  \item ( \quad ) Uso de objetos, classes e encapsulamento
  \item ( \quad ) Programação baseada em fatos e regras
  \item ( \quad ) Ênfase em funções puras e imutabilidade
  \item ( \quad ) Descrição do resultado esperado, não do algoritmo
  \item ( \quad ) Execução sequencial de comandos que alteram estado
\end{itemize}
\end{block}

\begin{block}{Questão 3 — Análise}
Considere o trecho de código abaixo:

\begin{verbatim}
const numeros = [1, 2, 3, 4, 5];
const dobro = numeros.map(n => n * 2);
\end{verbatim}

\begin{enumerate}
  \item Qual paradigma de programação está sendo utilizado?
  \item Cite duas características desse paradigma presentes no código.
  \item O vetor original foi modificado? Justifique.
\end{enumerate}
\end{block}

\begin{block}{Questão 4 — Paradigma Lógico}
Considere os fatos a seguir escritos em Prolog:

\begin{verbatim}
gosta(ana, cafe).
gosta(joao, cha).
gosta(ana, cha).
\end{verbatim}

Responda:
\begin{enumerate}
  \item O que representa cada fato?
  \item Qual seria a resposta para a consulta \texttt{?- gosta(ana, X).}?
  \item Como esse tipo de consulta difere de uma busca em linguagens imperativas?
\end{enumerate}
\end{block}

\begin{block}{Questão 5 — Comparação}
Explique duas diferenças entre:
\begin{itemize}
  \item Paradigma Orientado a Objetos e Paradigma Funcional;
  \item Paradigma Imperativo e Paradigma Declarativo.
\end{itemize}

Dê exemplos de situações em que cada paradigma seria mais adequado.
\end{block}

\begin{block}{Questão 6 — Reflexão}
Linguagens modernas como Java, Python e JavaScript são consideradas
\textbf{multiparadigma}.

Na sua opinião, isso é uma vantagem ou uma desvantagem?
Justifique sua resposta.
\end{block}

% =========================
\end{document}
% =========================
